Micropollutant records are full of \emph{non-detections}: the analyte was
sought and not found, so the result is a bound, not a value, yet stored as a
number. Article~3(3b) of Directive 2008/105/EC excludes such a result from
status assessment where its limit exceeds the standard, which a two-valued
record cannot express.

Auditing the European Environment Agency's Waterbase (\num{4190833} river
station-years, \num{637} substances, \num{37} countries) finds \num{44.9}\% of
samples below the quantification limit and \num{40.1}\% of assessments from a
method failing Directive 2009/90/EC. The two failures have different histories:
station-years declaring neither a censoring flag nor a limit --- \num{25.5}\%
overall --- are absent after \num{2012}, when the set of reporters also changes,
whereas since \num{2020} \num{19.3}\% of assessments are still below-quantification
results whose limit exceeds the standard, so the analytical failure is not an
artefact of an archive. Substitution manufactures findings: the same records yield
\num{59251} to \num{200708} exceedances as non-detections enter at zero, half
the limit or the limit, \num{112034} created by the substitution alone and only
\num{14809} affirmable in law. The \num{2026} standards are applied to a record that
predates them: the test is what the monitoring can decide, not who breached
what.

CENSO is an OWL~2 ontology on SOSA/SSN whose central choices the record decides.
The quantification limit is not a property of the instrument: \num{56.4}\% of
multi-year station--substance pairs carry more than one, which a threshold column
cannot hold, and CENSO carries the bound each limit implies onto the result it
qualifies. Compliance is three-valued --- met, not met, or not determinable, the
outcomes the Directive admits --- the third subtyped by the reason, which keeps
an unmeasurable standard apart from an unreported bound. Of \num{696168}
station-years a European standard reaches, \num{38.4}\% are undecidable, and the
largest single reason is analytical: \num{18.2}\% are below-quantification results
Article~3(3b) sets aside. The other half arises before any limit meets any
standard --- \num{19.2} points with no bound at all or with a standard written
over a quantity the record does not supply, once the water hardness the release
reports on rows of its own is joined and the first tier of the EU bioavailability
guidance is applied; read row by row, without that join, \num{43.8}\% are
undecidable. Swapping the regulation package re-decides
\num{16.0}\% of co-regulated outcomes. Of \num{21} ontologies parsed --- reaching outside the water domain to
the analytical-chemistry vocabularies, one of which defines the quantification
limit from the IUPAC Gold Book and still binds it to the method --- none carries
a censoring status, and \num{20} cannot separate a non-detection from a
measurement at the same number. In CENSO that distinction is an inconsistency,
not a label. Ontology, pipeline and figure data are open.
